\documentclass[11pt, twoside, twocolumn, paper = letter]{scrartcl}

\usepackage[autopunct]{csquotes}
\usepackage{fontspec}
\usepackage{microtype}
\usepackage{unicode-math}
\usepackage{lualatex-math}
\usepackage{hyperref}
\usepackage{xcolor}
\usepackage{luacolor}
\usepackage{scrlayer-scrpage}
\usepackage{tikz}

\usetikzlibrary{calc, math}

\setmainfont{STIX Two Text}
\setmathfont{STIX Two Math}
\setmonofont{Libertinus Mono}[Scale = MatchLowercase]
\setsansfont{Libertinus Sans}[Scale = MatchLowercase]

\usepackage{polyglossia}
\setdefaultlanguage[variant = american]{english}
\setotherlanguages{german}
\usepackage[
    sorting  = none
  , language = autobib
  , autolang = langname
  ]{biblatex}

\def\papertitle{Rational extrapolation for the Bulirsch--Stoer algorithm}

\addbibresource{rational-extrapolation.bib}

\hypersetup{
    unicode      = true
  , pdftoolbar   = true
  , pdfmenubar   = true
  , pdffitwindow = false
  , pdftitle     = {\papertitle}
  , pdfnewwindow = true
  , colorlinks   = true
  , linkcolor    = linkcolor
  , citecolor    = zzttqq
  , filecolor    = xdxdff
  , urlcolor     = uququq
}

\definecolor{linkcolor}{RGB}{98, 130, 41}
\definecolor{uququq}{rgb}{0.6,  0.2,  0}
\definecolor{zzttqq}{rgb}{0.6,  0.2,  0}
\definecolor{xdxdff}{rgb}{0.49, 0.49, 1}
\definecolor{titlecolor}{RGB}{69, 69, 73}
\definecolor{titlerulescolor}{RGB}{167, 184, 137}

\KOMAoptions{
    DIV      = 16
  , fontsize = 11
  , BCOR     = 8mm
  , pagesize = luatex
  }

\addtokomafont{title}{\color{titlecolor}}
\addtokomafont{pagehead}{\normalfont}

\title{
    \parbox[t]{\linewidth}{%
        \centering%
        \kern-\baselineskip%
        \kern-\ht\strutbox%
        {\color{titlerulescolor}\rule{\linewidth}{1pt}}\par%
        \kern-\dimexpr\baselineskip-\ht\strutbox\relax%
        \kern1ex%
        \papertitle\par%
        \kern1ex%
        \kern-.7\dimexpr\ht\strutbox+\dp\strutbox\relax%
        {\color{titlerulescolor}\rule{\linewidth}{1pt}}%
    }%
}
\date{}

\pagestyle{scrheadings}
\renewcommand*{\titlepagestyle}{scrheadings}

\begin{document}

\maketitle

\section{Rational extrapolation}

The base recurrence for rational extrapolation is given by Stoer and Bulirsch
\cite{BulirschStoer1964}. To avoid loss of significance from subtractive cancellation during
floating-point evaluation, we can apply the same difference technique they use for polynomial
extrapolation \cite{BulirschStoerBook}. A similar derivation can be found in \cite{NumRecipesInC}.

Let's assume we have $n + 1$ pairs of numbers $(x_i, f_i)$, $i = 0, \ldots, n$. We denote
$\tilde{T}_{ik}$ interpolation performed using pairs $i - k, i - k + 1, \ldots, i$. Neville's
algorithm for calculating rational interpolation $\tilde{T}_{nn}$ consists of applying these
recurrent equations:
\begin{equation*}
    \begin{gathered}
        \tilde{T}_{i0} = f_i
        \text{,}\quad
        \tilde{T}_{i,-1} = 0
        \text{,}\\
        \tilde{T}_{ik}
          = \tilde{T}_{i, k-1}
          + \frac{\tilde{T}_{i, k - 1} - \tilde{T}_{i - 1, k - 1}}
                 {
                     \frac{x - x_{i - k}}{x - x_i}
                     \left[
                         1- \frac{\tilde{T}_{i, k - 1} - \tilde{T}_{i - 1, k - 1}}
                                 {\tilde{T}_{i, k - 1} - \tilde{T}_{i - 1, k - 2}}
                     \right] - 1
                 }
        \text{.}
    \end{gathered}
\end{equation*}
For multistep integration we need $T_{ik} = \tilde{T}_{ik}(0)$, and also the rational function
depends on $h^2$ (rather than $h$), where $h$ is the current step. Therefore, the recurrent
equations take this form:
\begin{equation*}
    T_{ik}
      = T_{i, k-1}
      + \frac{T_{i, k - 1} - T_{i - 1, k - 1}}
             {
                 \left[\frac{h_{i - k}}{h_i}\right]^2
                 \left[
                     1- \frac{T_{i, k - 1} - T_{i - 1, k - 1}}
                             {T_{i, k - 1} - T_{i - 1, k - 2}}
                 \right] - 1
             }
    \text{.}
\end{equation*}
We can rewrite it as
\begin{equation*}
    T_{ik} = T_{i, k-1} + Q_{ik}
    \text{,}
\end{equation*}
by introducing the following small differences:
\begin{equation*}
    \begin{gathered}
        D_{ik} = T_{ik} - T_{i - 1, k - 1}
        \text{,}\\
        Q_{ik} = T_{ik} - T_{i, k - 1}
        \text{.}
    \end{gathered}
\end{equation*}
Using an obvious identity
\begin{equation*}
    D_{ik} - Q_{ik} = D_{i, k - 1} - Q_{i - 1, k - 1}
    \text{,}
\end{equation*}
we can deduce the recurrent equations for $D_{ik}$ and $Q_{ik}$:
\begin{equation*}
    \begin{gathered}
        Q_{ik}
          = D_{i, k - 1}
            \frac{D_{i, k - 1} - Q_{i - 1, k - 1}}
                 {
                    \left[\frac{h_{i - k}}{h_i}\right]^2 Q_{i - 1, k - 1} - D_{i, k - 1}
                 }
        \text{,}\\
        D_{ik}
          = \left[\frac{h_{i - k}}{h_i}\right]^2 Q_{i - 1, k - 1}
            \frac
            {
                D_{i, k - 1} - Q_{i - 1, k - 1}
            }
            {
                \left[\frac{h_{i - k}}{h_i}\right]^2 Q_{i - 1, k - 1} - D_{i, k - 1}
            }
        \text{.}
    \end{gathered}
\end{equation*}
Usually, $h_i$ have form $H / n_i$, where, for example, $n_i = 2, 4, 6, 8, 12, 16, \ldots$, $n_i = 2
n_{i - 2}$ for $i \geqslant 3$. Then the result doesn't depend on $H$, and we can write
\begin{equation*}
    \begin{gathered}
        Q_{i0} = D_{i0} = f_i
        \text{,}\\
        Q_{ik}
          = D_{i, k - 1}
            \frac{D_{i, k - 1} - Q_{i - 1, k - 1}}
                 {
                    \left[\frac{n_i}{n_{i - k}}\right]^2 Q_{i - 1, k - 1} - D_{i, k - 1}
                 }
        \text{,}\\
        D_{ik}
          = \left[\frac{n_i}{n_{i - k}}\right]^2 Q_{i - 1, k - 1}
            \frac
            {
                D_{i, k - 1} - Q_{i - 1, k - 1}
            }
            {
                \left[\frac{n_i}{n_{i - k}}\right]^2 Q_{i - 1, k - 1} - D_{i, k - 1}
            }
        \text{.}
    \end{gathered}
\end{equation*}
Finally,
\begin{equation*}
    T_{nn} = \sum_{k = 0}^n{Q_{nk}}
    \text{.}
\end{equation*}
To estimate an error, we introduce $U_{ik}$
\begin{equation*}
    U_{ik} = 2T_{i + 1, k} - T_{ik}
    \text{,}
\end{equation*}
which nears to the limit from the other side. Using definitions for $Q_{ik}$ and $D_{ik}$ we can
deduce
\begin{equation*}
    \begin{split}
        U_{i-1, k - 1} &= 2 T_{i, k - 1} - T_{i - 1, k - 1}\\
                       &= -2 Q_{ik} + 2 T_{ik} + D_{ik} - T_{ik}\\
                       &= T_{ik} + D_{ik} - 2 Q_{ik}\\
                       &= D_{ik} + T_{i - 1, k - 1} + D_{ik} - 2 Q_{ik}\\
                       &= T_{i - 1, k - 1} + 2\left(D_{ik} - Q_{ik}\right)\text{.}
    \end{split}
\end{equation*}
Therefore, an error estimate is
\begin{equation*}
    \varepsilon_{i - 1, k - 1}
      = \frac{1}{2}\left(T_{i-1, k - 1} - U_{i - 1, k - 1}\right) = Q_{ik} - D_{ik}
    \text{.}
\end{equation*}
Note that for the adaptive order and step control, Deuflhard suggested using $Q_{ik}$
\cite{Deuflhard1983}.

\section{Polynomial extrapolation for fallback}

The denominator in the recurrent expressions for the rational extrapolation can infinitely approach
zero. Stoer and Bulirsch in their book \cite{BulirschStoerBook} simply dismiss it by saying it never
happens in practice. However, it does happen in practice, and in this case, the entire integration
scheme blows up. To prevent this disaster from affecting the integration, we can switch the
corresponding $Q_{ik}$ and $D_{ik}$ to polynomial extrapolation. It works because every new row in
the extrapolation tableau is a refinement of the previous row. If we switch a refinement to a
different algorithm, it still will refine the previous values. For reference, below are the
corresponding recurrent relations for the polynomial extrapolation from \cite{BulirschStoerBook}.

Let's assume we have $n + 1$ pairs of numbers $(x_i, f_i)$, $i = 0, \ldots, n$. We denote
$\tilde{T}_{ik}$ interpolation performed using pairs $i - k, i - k + 1, \ldots, i$. Neville's
algorithm for calculating polynomial interpolation $\tilde{T}_{nn}$ consists of applying these
recurrent equations:
\begin{equation*}
    \begin{gathered}
        \tilde{T}_{i0} = f_i
        \text{,}\\
        \tilde{T}_{ik}
          = \tilde{T}_{i, k-1}
          + \frac{\tilde{T}_{i, k - 1} - \tilde{T}_{i - 1, k - 1}}
                 {
                     \frac{x - x_{i - k}}{x - x_i} - 1
                 }
        \text{.}
    \end{gathered}
\end{equation*}
Following the same derivation, we will get
\begin{equation*}
    \begin{aligned}
        Q_{i0} &= D_{i0} = f_i
        \text{,}\\
        Q_{ik}
         &= \frac{n_{i - k}^2}{n_i^2 - n_{i - k}^2}\left(D_{i, k - 1} - Q_{i - 1, k - 1}\right)
        \text{,}\\
        D_{ik}
         &= \frac{n_{i}^2}{n_i^2 - n_{i - k}^2}\left(D_{i, k - 1} - Q_{i - 1, k - 1}\right)
        \text{.}
    \end{aligned}
\end{equation*}
Similarly,
\begin{equation*}
    T_{nn} = \sum_{k = 0}^n{Q_{nk}}
    \text{.}
\end{equation*}

\printbibliography

\end{document}
